\setlength{\absparsep}{18pt} 

\begin{resumo}[Resumo]
 \begin{otherlanguage*}{brazil}
   
   \vspace{\onelineskip}
   \noindent 
   
O avanço de aplicações baseadas em inteligência artificial tornou comum que estudantes da Educação Básica convivam com sistemas inteligentes, sem que isso assegure a compreensão dos dados, regras e escolhas subjacentes ao seu funcionamento. Este projeto de qualificação investiga o ensino introdutório de \textit{machine learning} no Ensino Médio, concebido como conteúdo de Computação voltado a explicitar como modelos computacionais utilizam dados para produzir previsões, classificações e agrupamentos. O estudo tem como objetivo geral desenvolver e analisar uma sequência didática apoiada por uma plataforma interativa para o ensino de fundamentos de \textit{machine learning} a estudantes do Ensino Médio. A proposta pedagógica fundamenta-se nas normas curriculares da Resolução CNE/CEB nº 1/2022, no pensamento computacional e no aporte de múltiplas representações, estruturando-se no ciclo didático prever, experimentar, observar e explicar. O produto educacional articula uma sequência didática de oito encontros a uma plataforma web interativa executada localmente com dados sintéticos tabulares, contemplando cinco modelos introdutórios: regressão linear, regressão logística, \textit{k}-vizinhos mais próximos (\textit{k}-NN), árvore de decisão e \textit{k}-means. Metodologicamente, a pesquisa é de natureza aplicada, com abordagem predominantemente qualitativa e apoio de dados quantitativos descritivos, configurando-se como intervenção pedagógica realizada em uma turma de Ensino Médio sob docência do pesquisador. A produção de dados baseia-se em um diagnóstico inicial individual, registros docentes breves após cada aula e um desafio final individual de síntese com situações-problema inéditas. A análise qualitativa das produções examinará indícios de compreensão em quatro dimensões conceituais: reconhecimento de dados e finalidade da tarefa; antecipação de consequências de alterações; interpretação de representações; e discussão de limites de decisões automatizadas, subsidiando a sistematização do produto educacional e de parâmetros reflexivos para a atuação docente na escola pública.
   
   \textbf{Palavras-chave}: Ensino de Computação; Machine Learning; Ensino Médio; Sequência Didática; Plataforma Interativa; Produto Educacional.

 \end{otherlanguage*}
\end{resumo}

% O resumo em inglês (Abstract) será incluído após a validação da versão portuguesa.
